% ID: SEL-D0080-OA-E02
% Tanda S080.
\textbf{a)} Para una asíntota vertical en $x=-2$, debe anularse allí el denominador: $-2+b=0$, de donde $b=2$. El límite en ambos infinitos de $ax/(x+b)$ es $a$, por lo que la horizontal $y=3$ exige $a=3$. Con estos valores, el numerador en $-2$ es $-6$, que no se anula, confirmando la asíntota vertical. Por tanto, $\boxed{a=3,\ b=2}$.

\textbf{b)} El dominio del polinomio es $\mathbb R$. Su derivada es $g'(x)=3x(x-2)$, positiva para $x<0$ y $x>2$, y negativa para $0<x<2$. Así, crece en $(-\infty,0)$ y $(2,\infty)$ y decrece en $(0,2)$. Los cambios de signo dan un máximo relativo en $(0,2)$ y un mínimo relativo en $(2,-2)$.

El esbozo pasa por ambos extremos y va de $-\infty$ a $+\infty$ al recorrer el eje de izquierda a derecha.
\begin{center}
\setlength{\unitlength}{1mm}
\begin{picture}(104,70)(-28,-35)
\put(-24,0){\vector(1,0){94}}\put(0,-29){\vector(0,1){59}}
\put(71,0){$x$}\put(1,30){$g$}
\qbezier(-20.0000,-20.0000)(-18.0000,-11.0200)(-16.0000,-4.3200)
\qbezier(-16.0000,-4.3200)(-14.0000,2.3800)(-12.0000,7.0400)
\qbezier(-12.0000,7.0400)(-10.0000,11.7000)(-8.0000,14.5600)
\qbezier(-8.0000,14.5600)(-6.0000,17.4200)(-4.0000,18.7200)
\qbezier(-4.0000,18.7200)(-2.0000,20.0200)(0.0000,20.0000)
\qbezier(0.0000,20.0000)(2.0000,19.9800)(4.0000,18.8800)
\qbezier(4.0000,18.8800)(6.0000,17.7800)(8.0000,15.8400)
\qbezier(8.0000,15.8400)(10.0000,13.9000)(12.0000,11.3600)
\qbezier(12.0000,11.3600)(14.0000,8.8200)(16.0000,5.9200)
\qbezier(16.0000,5.9200)(18.0000,3.0200)(20.0000,0.0000)
\qbezier(20.0000,0.0000)(22.0000,-3.0200)(24.0000,-5.9200)
\qbezier(24.0000,-5.9200)(26.0000,-8.8200)(28.0000,-11.3600)
\qbezier(28.0000,-11.3600)(30.0000,-13.9000)(32.0000,-15.8400)
\qbezier(32.0000,-15.8400)(34.0000,-17.7800)(36.0000,-18.8800)
\qbezier(36.0000,-18.8800)(38.0000,-19.9800)(40.0000,-20.0000)
\qbezier(40.0000,-20.0000)(42.0000,-20.0200)(44.0000,-18.7200)
\qbezier(44.0000,-18.7200)(46.0000,-17.4200)(48.0000,-14.5600)
\qbezier(48.0000,-14.5600)(50.0000,-11.7000)(52.0000,-7.0400)
\qbezier(52.0000,-7.0400)(54.0000,-2.3800)(56.0000,4.3200)
\qbezier(56.0000,4.3200)(58.0000,11.0200)(60.0000,20.0000)
\put(2,23){\scriptsize $(0,2)$}
\put(34,-27){\scriptsize $(2,-2)$}
\end{picture}
\end{center}
